\section*{Anexo V. Autodeclaración del uso de IA}
Se declara que durante el desarrollo de este proyecto se han utilizado herramientas basadas en Inteligencia Artificial generativa como apoyo técnico, operativo y de consulta. El uso de dichas tecnologías se ha realizado bajo la constante supervisión, revisión y criterio del autor, garantizando en todo momento la integridad académica del trabajo.

Concretamente, la asistencia de la IA se ha aplicado en las siguientes tareas:
\begin{itemize}
    \item \textbf{Prototipado y estructuración inicial:} Generación de la estructura base del proyecto y desarrollo de un frontend provisional, lo que permitió agilizar las fases iniciales de diseño y visualización de la interfaz.
    \item \textbf{Soporte en codificación y refactorización:} Asistencia técnica en el proceso de migración de módulos del frontend hacia JavaScript, así como apoyo en la resolución de bloqueos lógicos, optimización de sintaxis y tareas de debugging.
    \item \textbf{Generación de pruebas:} Apoyo en la redacción y estructuración de casos de prueba unitarios para diversas funciones, contribuyendo a asegurar la robustez y calidad del software desarrollado.
    \item \textbf{Documentación técnica:} Asistencia en la generación de comentarios estructurados y documentación de métodos y funciones, mejorando la legibilidad del código fuente.
    \item \textbf{Consulta y síntesis de información:} Uso de la herramienta como asistente de consulta técnica para consultar rápidamente conceptos específicos y resolver dudas de implementación, actuando como complemento y síntesis de la documentación oficial de las bibliotecas utilizadas.
    \item \textbf{Maquetación y conversión a LaTeX:} Transformación de documentos de Google Docs al lenguaje de marcado \LaTeX, incluyendo la automatización del formato de tablas y la organización del glosario de términos.
    \item \textbf{Revisión de estilo:} Apoyo en la corrección gramatical y el pulido tipográfico del documento final para asegurar un tono académico coherente.
\end{itemize}

Es importante destacar que el diseño de la arquitectura del software, la toma de decisiones tecnológicas, la implementación de la lógica de negocio central, así
 como la interpretación de los resultados, son fruto exclusivo del trabajo del autor. La inteligencia artificial ha actuado únicamente
  como una herramienta de productividad y soporte técnico, manteniendo intacta la autoría y la responsabilidad sobre el contenido total de este Trabajo de Fin de Grado.